\section{PageRank}

\subsection{PageRank}

Voici un graphe repr�sentant les hyperliens entre les pages d'un site web :\\

\fig{RI/PageRank1.eps}{Graphe d'hyperliens entre des pages}{6cm}

\begin{enumerate}
	\item Rappeler la loi de probabilit� de passer d'une page $i$ � une page $j$.
	\begin{correction}
		P$_{ij} = \frac{\alpha}{N} + \frac{I (1 - \alpha)}{K}$\\
		O� $N$ est le nombre de pages, $K$ le nombre d'hyperliens de la page $i$.\\
		I est �gal 1 si $j$ est une des cibles des hyperliens, 0 sinon.
	\end{correction}

	\item Construire la matrice NxN de probabilit�s de passages d'une page � une autre, avec $\alpha = \frac{1}{2}$
	\begin{correction}Nous avons N=4, K$_A$=2, K$_B$=2, K$_C$=2, K$_D$=1\\
		\begin{tabular}{|c|c|c|c|}
		\hline
		1/8 & 3/8 & 1/8 & 3/8\\
		\hline
		1/8 & 1/8 & 3/8 & 3/8\\
		\hline
		1/8 & 3/8 & 1/8 & 3/8\\
		\hline
		5/8 & 1/8 & 1/8 & 1/8\\
		\hline
		\end{tabular}
	\end{correction}

	\item Calculer le PageRank de la matrice � l'aide du vecteur (1 0 0 0). Vous arrondirez chaque calcul avec 3 chiffres apr�s la virgule.
	\begin{correction}Etapes :\\
\begin{tabular}{|c|llll|}
\hline
1&   1        &0        &0        &0 \\
\hline
1&		0,125&	0,375&	0,125&	0,375\\
\hline
2&		0,313&	0,188&	0,219&	0,281\\
\hline
3&		0,266&	0,258&	0,172&	0,305\\
\hline
4&		0,278&	0,235&	0,190&	0,299\\
\hline
5&		0,275&	0,242&	0,184&	0,301\\
\hline
6&		0,276&	0,240&	0,186&	0,301\\
\hline
7&		0,276&	0,241&	0,185&	0,301\\
\hline
8&		0,276&	0,241&	0,186&	0,301\\
\hline
9&		0,276&	0,241&	0,186&	0,301\\
\hline
\end{tabular}
	\end{correction}
\end{enumerate}

\subsection{Pertinence du contenu d'une page Web}

\begin{enumerate}
\item En reprenant le r�sultat de la recherche par cosinus (section~\ref{cos}, question~\ref{cos_q}) sur les documents $doc_1$ (site A), $doc_2$ (site B), $doc_3$ (site C), $doc_4$ (site D), et en multipliant le score de cosinus par le pageRank du site h�bergeant ce texte, donner l'ordre final de votre recherche.
\begin{correction}
Nous obtenons ainsi un PageRank pour chaque page de :
\begin{itemize}
\item $doc_1$ : 0,276 $\times$ 0,126 = 0,035
\item $doc_3$ : 0,186 $\times$ 0,094 = 0,018
\item $doc_2$ : 0,241 $\times$ 0,036 = 0,009
\item $doc_4$ : 0,301 $\times$ 0,021 = 0,006
\end{itemize}
\end{correction}

\end{enumerate}
